% Option C: Organize by model concepts (lifecycle events / effects / coordination)
% Only the Policy Interfaces + Verification sections are restructured.
% Design Principles, Challenges, and High-Level Architecture are unchanged.
% All code examples and technical content preserved from original.

% ============================================================
% REPLACES: §§ Policy Interfaces + §§ Runtime Verification
% ============================================================

\subsection{Lifecycle Events and Triggers}
\label{sec:design-interfaces}

GPU resources follow driver-defined lifecycles. Memory chunks transition between activated, accessed, and evicted; compute contexts transition between created, scheduled, and preempted. \sys exposes these lifecycle transitions as programmable events. BPF programs attach at events across three layers, each providing context that the others cannot.

\subsubsection{Driver-Level Events (struct\_ops)}

The host driver exports memory and scheduling lifecycle events through \texttt{struct\_ops} tables. Memory events include four transitions (\emph{activate}, \emph{access}, \emph{evict\_prepare}, and \emph{prefetch}). Scheduling events cover hardware-queue lifecycle (\emph{task\_init}, \emph{task\_destroy}). Because memory and scheduling share the same fault handler, a single BPF program can observe both domains simultaneously.

\sys regions correspond to GPU hardware-managed units, such as 2MB physical UVM chunks or VA blocks, and finer-grained 4KB pages. This granularity balances policy precision against metadata overhead. Using the regions abstraction, \sys hides vendor-specific and hardware complexities. Handlers encode decisions through an enum field in the context and may reorder but never remove regions in the kernel-maintained eviction list. The kernel enforces safety and correctness, including fallback FIFO eviction under pressure.

\begin{lstlisting}
// Host-side lifecycle events: memory
struct gpu_mem_ops {
  int (*gpu_activate)(gdrv_mem_add_ctx_t *ctx);
  int (*gpu_access)(gdrv_mem_access_ctx_t *ctx);
  int (*gpu_evict_prepare)(gdrv_mem_remove_ctx_t *ctx);
  int (*gpu_prefetch)(gdrv_mem_prefetch_ctx_t *ctx);
};

// Host-side lifecycle events: scheduling
struct gpu_sched_ops {
  int  (*task_init)(gsched_queue_ctx_t *ctx);
  void (*task_destroy)(gsched_queue_ctx_t *ctx);
};
\end{lstlisting}

\subsubsection{Application-Level Triggers (uprobes)}

\sys attaches uprobes to unmodified CUDA API calls (e.g., \texttt{cudaLaunchKernel}, \texttt{cudaStreamSynchronize}), enabling proactive policy decisions from application-level events that are invisible to the driver. These triggers are decoupled from lifecycle transitions: a BPF program at \texttt{cudaLaunchKernel} can direct future memory transitions (prefetch) and scheduling transitions (priority) before any fault occurs.

\subsubsection{Device-Level Events (device-side BPF)}

Device policies are triggered at memory operations and thread-block scheduling events on the GPU, receiving a compact warp-uniform context. These events expose GPU-internal execution state that is architecturally opaque to the host.

\begin{lstlisting}
// Device-side events: memory
struct gdev_mem_ops {
    void (*access)(gdev_mem_access_ctx_t *ctx);
    void (*fence)(gdev_mem_fence_ctx_t *ctx);
};

// Device-side events: scheduling
struct gdev_sched_ops {
  void (*enter)(gdev_block_ctx_t *ctx);
  void (*exit)(gdev_block_ctx_t *ctx);
  void (*probe)(gdev_block_ctx_t *ctx);
  void (*retprobe)(gdev_block_ctx_t *ctx);
  bool (*should_try_steal)(gdev_block_ctx_t *ctx);
};
\end{lstlisting}

% ============================================================

\subsection{Effects}
\label{sec:design-effects}

BPF programs direct lifecycle transitions via effects. All effects are delivered through kfuncs, kernel functions exposed by the driver. Effects have different timescales and scopes:

\subsubsection{Synchronous Effects (Host)}

Host-side kfuncs reorder the eviction list, set scheduling priority and timeslice, or reject queue creation. These take effect immediately within the current hook invocation.

\begin{lstlisting}
// Memory effects
void bpf_gpu_move_head(gdrv_mem_list_t *list,
                               gmem_region_t *r);
void bpf_gpu_move_tail(gdrv_mem_list_t *list,
                               gmem_region_t *r);

// Scheduling effects
void bpf_gpu_set_attr(gsched_queue_ctx_t *ctx, u64 us);
void bpf_gpu_reject_bind(gsched_queue_ctx_t *ctx);
\end{lstlisting}

\subsubsection{Asynchronous Effects (Cross-Device)}

Because cross-device transitions take milliseconds, \sys provides asynchronous kfuncs that initiate effects after the synchronous hook returns. BPF programs schedule deferred work via \texttt{bpf\_wq}, enabling cross-block prefetch, proactive migration, and cooperative preemption through the driver's native context-switch mechanism. Policies set queue priority, timeslice, and may reject or defer queue creation. Properties are written directly into firmware-visible structures, ensuring enforcement beyond user-space hints.

\subsubsection{Device-Side Effects}

Device-side helpers enable on-device policy effects. Operations like prefetch can be performed on device and then trigger host-side prefetch handlers, enabling the host to combine driver information with device-side predictive logic. The work-stealing thread-block scheduler allows device-side eBPF handlers to steer block scheduling decisions. GPU kernels expose their logical work units; persistent worker blocks select and process units. The verifier ensures that scheduling decisions respect GPU kernel invariants. \sys can inject this scheduler via binary rewriting or allow developers to integrate it.

\begin{lstlisting}
// Device-side helpers
int gdev_mem_prefetch(gmem_region_t* r);
\end{lstlisting}

% ============================================================

\subsection{Cross-Device State Coordination}

\subsubsection{Hierarchical Cross-layer Maps}

GPU policy decisions require coherent sharing of policy state across asynchronous CPU and GPU contexts. To transparently bridge CPU-GPU memory hierarchies, \sys introduces cross-layer maps. Logically, each map is a single unified key-value store accessible from host-side driver hooks, GPU-side device policies, and user-space control planes, and verifiable through eBPF's existing model. Physically, maps are realized as hierarchical shards across host DRAM, GPU global memory, and GPU-local storage (e.g., shared memory, SM-local). The runtime dynamically determines shard placement based on access frequency, latency, and capacity constraints. Following eBPF's soft-state philosophy, \sys maps provide relaxed, eventual consistency. Instead of strong global ordering, our runtime periodically merges GPU-local map shards into canonical snapshots at synchronization points, such as GPU kernel completion boundaries. This snapshot-based model is sufficient for policy decisions relying on approximate statistics. Occasional staleness affects decision optimality but cannot violate correctness invariants such as memory integrity, which remain enforced by the GPU driver and hardware MMU.

% ============================================================

\subsection{Verification and Safety}
\label{sec:verification}

\sys policies execute directly within privileged GPU driver and kernel contexts, making safety guarantees critical. We adopt a standard eBPF threat model: we trust system administrators loading policies, but consider policy code potentially buggy or misconfigured. The trusted computing base (TCB) includes the kernel and \sys kernel module, GPU compiler backend, and GPU driver/firmware. Policy authors only interact with eBPF and do not write native device code. For host-side verification, we reuse the existing Linux kernel eBPF verifier unchanged, relying on standard type-checking, bounded-loop, and memory-safety constraints. \sys-specific kfuncs and hooks are declared via BPF Type Format (BTF) metadata.

\subsubsection{SIMT-aware Verification}

\sys introduces a SIMT-aware static verification model enforcing warp-level uniformity constraints to ensure efficient and safe GPU-side policy execution.
The key insight is that we distinguish \emph{warp-uniform} values (identical across warp threads, e.g., region identifiers) from \emph{lane-varying} values (thread-local GPU state). Policies may compute freely using lane-varying values, but these must not influence control-flow decisions or directly produce side effects visible outside a warp. To guarantee warp-level coherence, \sys mandates that branch conditions, loop bounds, and map-update keys be warp-uniform or derived via warp-level aggregation. Additionally, \sys prohibits GPU-wide barriers, global synchronization primitives, and non-uniform atomic operations, and imposes resource budgets per policy hook to bound memory and thread resource usage. These constraints ensure safe policy integration that respects GPU hardware execution semantics.

\subsubsection{Warp-level Execution Optimizations}

Direct scalar execution of eBPF after SIMT verification on GPU threads still causes resource overhead (e.g. duplicate memory update to maps) due to the SIMT. To resolve this, \sys introduces a warp-level aggregated execution model, executing policy logic once per warp using a designated "warp leader" (typically lane 0). Each hook invocation conceptually follows a two-phase approach: each lane independently computes lane-local contributions (e.g., effective addresses, counters) without influencing control flow or directly updating shared state. These contributions are then aggregated at warp granularity. Subsequently, the warp leader executes the policy handler once using aggregated inputs and warp-uniform context. The resulting decision is broadcast back to all lanes. This design minimizes GPU overhead while preserving eBPF's scalar semantics. To further reduce runtime overhead, \sys employs compiler-level optimizations (specialization and inlining) tailored for warp-level execution.
